\documentclass[14pt]{extarticle}

\usepackage{extsizes}
\usepackage[T1,T2A]{fontenc}
\usepackage[utf8]{inputenc}
\pagestyle{plain}

\usepackage[a4paper,width=170mm,top=20mm,bottom=20mm]{geometry}

\usepackage{amsmath}
\usepackage{amssymb}

\makeatletter
\renewcommand\thesection{}
\renewcommand\thesubsection{}
\makeatother

\usepackage{hyperref}
\hypersetup{
	colorlinks=true,
	linkcolor=blue,
	citecolor=blue,
	urlcolor=magenta
}


\begin{document}
	\section*{Контрольная работа}
	Цель контрольной работы: определение уровня [остаточных] знаний по курсу компьютерной алгебры (1-й семестр).

	Каждому студенту предназначен индивидуальный вариант, номер которого совпадает с номером студента в списке группы. Вариант состоит из семи заданий.

	\textbf{Обратите внимание}: варианты для студентов группы n приведены в разделе «Варианты группы n»!

	Оценка $s$ вычисляется по формуле
	\[
	s = \begin{cases}
		0,&\textrm{s < 3,}\\
		3,&\textrm{3 $\leq$ s < 5,}\\
		4,&\textrm{5 $\leq$ s < 6,}\\
		5,&\textrm{s $\geq$ 6.}
		\end{cases},
	\]
	где $b$ --- количество набранных баллов за основные задания.

	Критерии оценивания:
	\begin{itemize}
		\item если при выполнении задания студент верно и точно применил требуемый алгоритм, полно и понятно описал шаги и действия и получил верный ответ, задание оценивается одним баллом;
		\item если студент верно и точно применил требуемый алгоритм, полно и понятно описал шаги и действия, но получил неверный ответ из-за допущенной вычислительной ошибки, задание оценивается половиной балла;
		\item если студент верно и точно применил требуемый алгоритм, но не полно или не понятно описал шаги и действия, но получил верный ответ, задание также оценивается половиной балла;
		\item в случаях, когда решение не удовлетворяет изложенным выше требованиям, задание оценивается в 0 баллов.
	\end{itemize}

	Оценка $5$ за КР карается "полуавтоматом" на зачёте (минус один вопрос).

	При решении КР можно пользоваться всем, чем угодно, кроме помощи одногруппников.

	\clearpage

    \tableofcontents 

	\clearpage
	\section{Варианты группы 6411}
	\subsection{Вариант 1}
	\begin{enumerate}
		\item Вычислить произведение целых чисел \[33203408 \cdot 1340221\] с помощью метода Карацубы. Применять метод рекурсивно, база рекурсии — двузначные числа, которые можно умножать любым способом.\\ Использовать десятичную систему счисления.
		\item Вычислить произведение целых чисел \[661 \cdot 115\] с помощью БПФ в поле вычетов $\mathbb{Z}/p\mathbb{Z}$. Обосновать выбор параметров алгоритма.
		\item Вычислить произведение в поле Ферма: \[35 \cdot 80\ (\textrm{mod}\ 257).\]
		\item Возвести целое число в степень по модулю с помощью бинарного алгоритма: \[838^{54}\ (\textrm{mod}\ 415).\]
		\item С помощью теста Миллера-Рабина определить, является ли простым число $N = 113$. Требуется дать ответ с вероятностью ошибки не~более~$\frac{1}{16}$.
		\item Определить порядок элемента $205$ поля вычетов $\mathbb{Z}/463\mathbb{Z}$ с помощью построения дерева произведений.
		\item Открытый текст $M$ зашифрован криптосистемой RSA с помощью пары открытых ключей $(e_1, N)$ и $(e_2, N)$:  \[\begin{aligned} C_1 \equiv M^{e_1}\ (\textrm{mod}\ N), \\ C_2 \equiv M^{e_2}\ (\textrm{mod}\ N). \end{aligned}\] Зная пару шифротекстов и открытые ключи: \[C_1 = 102541, C_2 = 47172, e_1 = 11, e_2 = 13, N = 109999,\] восстановить открытый текст $M$ без разложения модуля $N$ на множители.
	\end{enumerate}
\clearpage	\subsection{Вариант 2}
	\begin{enumerate}
		\item Вычислить произведение целых чисел по модулю: \[157733 \cdot 929051\ (\textrm{mod}\ 6184140)\] с помощью метода Монтгомери. Обосновать выбор параметра $R$.
		\item Вычислить произведение целых чисел \[58 \cdot 901\] с помощью БПФ в поле вычетов $\mathbb{Z}/p\mathbb{Z}$. Обосновать выбор параметров алгоритма.
		\item Вычислить произведение в поле Мерсенна: \[188 \cdot 101\ (\textrm{mod}\ 255).\]
		\item Возвести целое число в степень с помощью бинарного алгоритма: \[3^{101}.\]
		\item С помощью теста Миллера-Рабина определить, является ли простым число $N = 881$. Требуется дать ответ с вероятностью ошибки не~более~$\frac{1}{16}$.
		\item Определить порядок элемента $206$ поля вычетов $\mathbb{Z}/331\mathbb{Z}$ с помощью построения дерева произведений.
		\item Открытый текст $M$ зашифрован криптосистемой RSA с помощью пары открытых ключей $(e_1, N)$ и $(e_2, N)$:  \[\begin{aligned} C_1 \equiv M^{e_1}\ (\textrm{mod}\ N), \\ C_2 \equiv M^{e_2}\ (\textrm{mod}\ N). \end{aligned}\] Зная пару шифротекстов и открытые ключи: \[C_1 = 48875, C_2 = 60480, e_1 = 29, e_2 = 23, N = 64517,\] восстановить открытый текст $M$ без разложения модуля $N$ на множители.
	\end{enumerate}
\clearpage	\subsection{Вариант 3}
	\begin{enumerate}
		\item Вычислить произведение целых чисел по модулю: \[613473 \cdot 245432\ (\textrm{mod}\ 3164242)\] с помощью метода Монтгомери. Обосновать выбор параметра $R$.
		\item Вычислить произведение целых чисел \[870 \cdot 83\] с помощью БПФ в поле вычетов $\mathbb{Z}/p\mathbb{Z}$. Обосновать выбор параметров алгоритма.
		\item Вычислить произведение в поле Мерсенна: \[253 \cdot 187\ (\textrm{mod}\ 255).\]
		\item Возвести целое число в степень с помощью $4$-арного алгоритма: \[2^{121}.\]
		\item С помощью теста Миллера-Рабина определить, является ли простым число $N = 881$. Требуется дать ответ с вероятностью ошибки не~более~$\frac{1}{16}$.
		\item Определить порядок элемента $469$ поля вычетов $\mathbb{Z}/967\mathbb{Z}$ с помощью построения дерева произведений.
		\item Восстановить параметры $a$, $c$ и $m$ Линейного конгруэнтного генератора \[x_i \equiv a \cdot x_{i-1} + c\ (\textrm{mod}\ m)\] по последовательности подряд идущих выходных значений генератора: \[129, 147, 31, 313, 32, 260.\]
	\end{enumerate}
\clearpage	\subsection{Вариант 4}
	\begin{enumerate}
		\item Вычислить произведение целых чисел по модулю: \[318674 \cdot 882250\ (\textrm{mod}\ 5513676)\] с помощью метода Монтгомери. Обосновать выбор параметра $R$.
		\item Вычислить произведение целых чисел \[42 \cdot 66\] с помощью БПФ в поле вычетов $\mathbb{Z}/p\mathbb{Z}$. Обосновать выбор параметров алгоритма.
		\item Вычислить произведение в поле Мерсенна: \[45 \cdot 174\ (\textrm{mod}\ 255).\]
		\item Возвести целое число в степень по модулю с помощью $4$-арного алгоритма: \[905^{118}\ (\textrm{mod}\ 417).\]
		\item С помощью теста Миллера-Рабина определить, является ли простым число $N = 977$. Требуется дать ответ с вероятностью ошибки не~более~$\frac{1}{16}$.
		\item Определить порядок элемента $499$ поля вычетов $\mathbb{Z}/547\mathbb{Z}$ с помощью построения дерева произведений.
		\item Восстановить параметры $a$, $c$ и $m$ Линейного конгруэнтного генератора \[x_i \equiv a \cdot x_{i-1} + c\ (\textrm{mod}\ m)\] по последовательности подряд идущих выходных значений генератора: \[255, 127, 129, 244, 108, 59.\]
	\end{enumerate}
\clearpage	\subsection{Вариант 5}
	\begin{enumerate}
		\item Вычислить произведение целых чисел \[750165815 \cdot 779355825\] с помощью метода Тоома-Кука. Использовать полином 2-й степени и десятичную систему счисления.
		\item Вычислить произведение целых чисел \[836 \cdot 560\] с помощью БПФ в поле вычетов $\mathbb{Z}/p\mathbb{Z}$. Обосновать выбор параметров алгоритма.
		\item Вычислить произведение в поле Мерсенна: \[73 \cdot 122\ (\textrm{mod}\ 255).\]
		\item Возвести целое число в степень с помощью $4$-арного алгоритма: \[2^{96}.\]
		\item С помощью теста Миллера-Рабина определить, является ли простым число $N = 953$. Требуется дать ответ с вероятностью ошибки не~более~$\frac{1}{16}$.
		\item Определить порядок элемента $138$ поля вычетов $\mathbb{Z}/331\mathbb{Z}$ с помощью построения дерева произведений.
		\item Открытый текст $M$ зашифрован криптосистемой RSA с помощью пары открытых ключей $(e_1, N)$ и $(e_2, N)$:  \[\begin{aligned} C_1 \equiv M^{e_1}\ (\textrm{mod}\ N), \\ C_2 \equiv M^{e_2}\ (\textrm{mod}\ N). \end{aligned}\] Зная пару шифротекстов и открытые ключи: \[C_1 = 23241, C_2 = 12115, e_1 = 19, e_2 = 13, N = 26989,\] восстановить открытый текст $M$ без разложения модуля $N$ на множители.
	\end{enumerate}
\clearpage
	\section{Варианты группы 6412}
	\subsection{Вариант 1}
	\begin{enumerate}
		\item Вычислить произведение целых чисел \[633457448 \cdot 596897114\] с помощью метода Тоома-Кука. Использовать полином 2-й степени и десятичную систему счисления.
		\item Вычислить произведение целых чисел \[987 \cdot 889\] с помощью БПФ в поле вычетов $\mathbb{Z}/p\mathbb{Z}$. Обосновать выбор параметров алгоритма.
		\item Вычислить произведение в поле Мерсенна: \[28 \cdot 163\ (\textrm{mod}\ 255).\]
		\item Возвести целое число в степень с помощью бинарного алгоритма: \[2^{73}.\]
		\item С помощью теста Соловея-Штрассена определить, является ли простым число $N = 401$. Требуется дать ответ с вероятностью ошибки не~более~$\frac14$.
		\item Определить порядок элемента $47$ поля вычетов $\mathbb{Z}/859\mathbb{Z}$ с помощью построения дерева произведений.
		\item Открытый текст $M$ зашифрован криптосистемой RSA с помощью пары открытых ключей $(e_1, N)$ и $(e_2, N)$:  \[\begin{aligned} C_1 \equiv M^{e_1}\ (\textrm{mod}\ N), \\ C_2 \equiv M^{e_2}\ (\textrm{mod}\ N). \end{aligned}\] Зная пару шифротекстов и открытые ключи: \[C_1 = 16744, C_2 = 7891, e_1 = 23, e_2 = 11, N = 54841,\] восстановить открытый текст $M$ без разложения модуля $N$ на множители.
	\end{enumerate}
\clearpage	\subsection{Вариант 2}
	\begin{enumerate}
		\item Вычислить произведение целых чисел \[652244083 \cdot 950591175\] с помощью метода Тоома-Кука. Использовать полином 2-й степени и десятичную систему счисления.
		\item Вычислить произведение целых чисел \[742 \cdot 549\] с помощью БПФ в поле вычетов $\mathbb{Z}/p\mathbb{Z}$. Обосновать выбор параметров алгоритма.
		\item Вычислить произведение в поле Ферма: \[81 \cdot 28\ (\textrm{mod}\ 257).\]
		\item Возвести целое число в степень по модулю: \[28^{127}\ (\textrm{mod}\ 340)\] с помощью метода Монтгомери. Обосновать выбор параметра $R$.
		\item С помощью теста Соловея-Штрассена определить, является ли простым число $N = 409$. Требуется дать ответ с вероятностью ошибки не~более~$\frac14$.
		\item Определить порядок элемента $315$ поля вычетов $\mathbb{Z}/859\mathbb{Z}$ с помощью построения дерева произведений.
		\item Восстановить параметры $a$, $c$ и $m$ Линейного конгруэнтного генератора \[x_i \equiv a \cdot x_{i-1} + c\ (\textrm{mod}\ m)\] по последовательности подряд идущих выходных значений генератора: \[414, 193, 96, 262, 168, 89.\]
	\end{enumerate}
\clearpage	\subsection{Вариант 3}
	\begin{enumerate}
		\item Вычислить произведение целых чисел \[7895486 \cdot 47575947\] с помощью метода Карацубы. Применять метод рекурсивно, база рекурсии — двузначные числа, которые можно умножать любым способом.\\ Использовать десятичную систему счисления.
		\item Вычислить произведение целых чисел \[694 \cdot 394\] с помощью БПФ в поле вычетов $\mathbb{Z}/p\mathbb{Z}$. Обосновать выбор параметров алгоритма.
		\item Вычислить произведение в поле Мерсенна: \[246 \cdot 204\ (\textrm{mod}\ 255).\]
		\item Возвести целое число в степень с помощью бинарного алгоритма: \[2^{45}.\]
		\item С помощью теста Соловея-Штрассена определить, является ли простым число $N = 479$. Требуется дать ответ с вероятностью ошибки не~более~$\frac14$.
		\item Определить порядок элемента $410$ поля вычетов $\mathbb{Z}/691\mathbb{Z}$ с помощью построения дерева произведений.
		\item Восстановить параметры $a$, $c$ и $m$ Линейного конгруэнтного генератора \[x_i \equiv a \cdot x_{i-1} + c\ (\textrm{mod}\ m)\] по последовательности подряд идущих выходных значений генератора: \[23, 2, 38, 4, 90, 95.\]
	\end{enumerate}
\clearpage	\subsection{Вариант 4}
	\begin{enumerate}
		\item Вычислить произведение целых чисел \[87657988 \cdot 72126881\] с помощью метода Карацубы. Применять метод рекурсивно, база рекурсии — двузначные числа, которые можно умножать любым способом.\\ Использовать десятичную систему счисления.
		\item Вычислить произведение целых чисел \[256 \cdot 837\] с помощью БПФ в поле вычетов $\mathbb{Z}/p\mathbb{Z}$. Обосновать выбор параметров алгоритма.
		\item Вычислить произведение в поле Мерсенна: \[232 \cdot 139\ (\textrm{mod}\ 255).\]
		\item Возвести целое число в степень по модулю с помощью бинарного алгоритма: \[918^{59}\ (\textrm{mod}\ 341).\]
		\item С помощью теста Миллера-Рабина определить, является ли простым число $N = 409$. Требуется дать ответ с вероятностью ошибки не~более~$\frac{1}{16}$.
		\item Определить порядок элемента $314$ поля вычетов $\mathbb{Z}/631\mathbb{Z}$ с помощью построения дерева произведений.
		\item Восстановить параметры $a$, $c$ и $m$ Линейного конгруэнтного генератора \[x_i \equiv a \cdot x_{i-1} + c\ (\textrm{mod}\ m)\] по последовательности подряд идущих выходных значений генератора: \[68, 114, 155, 137, 79, 128.\]
	\end{enumerate}
\clearpage	\subsection{Вариант 5}
	\begin{enumerate}
		\item Вычислить произведение целых чисел \[567631829 \cdot 527600234\] с помощью метода Тоома-Кука. Использовать полином 2-й степени и десятичную систему счисления.
		\item Вычислить произведение целых чисел \[488 \cdot 155\] с помощью БПФ в поле вычетов $\mathbb{Z}/p\mathbb{Z}$. Обосновать выбор параметров алгоритма.
		\item Вычислить произведение в поле Мерсенна: \[141 \cdot 210\ (\textrm{mod}\ 255).\]
		\item Возвести целое число в степень с помощью $4$-арного алгоритма: \[3^{54}.\]
		\item С помощью теста Миллера-Рабина определить, является ли простым число $N = 449$. Требуется дать ответ с вероятностью ошибки не~более~$\frac{1}{16}$.
		\item Определить порядок элемента $4$ поля вычетов $\mathbb{Z}/463\mathbb{Z}$ с помощью построения дерева произведений.
		\item Открытый текст $M$ зашифрован криптосистемой RSA с помощью пары открытых ключей $(e_1, N)$ и $(e_2, N)$:  \[\begin{aligned} C_1 \equiv M^{e_1}\ (\textrm{mod}\ N), \\ C_2 \equiv M^{e_2}\ (\textrm{mod}\ N). \end{aligned}\] Зная пару шифротекстов и открытые ключи: \[C_1 = 44262, C_2 = 111178, e_1 = 3, e_2 = 19, N = 112669,\] восстановить открытый текст $M$ без разложения модуля $N$ на множители.
	\end{enumerate}
\clearpage
	\section{Варианты группы 6413}
	\subsection{Вариант 1}
	\begin{enumerate}
		\item Вычислить произведение целых чисел по модулю: \[125022 \cdot 660940\ (\textrm{mod}\ 5705544)\] с помощью метода Монтгомери. Обосновать выбор параметра $R$.
		\item Вычислить произведение целых чисел \[602 \cdot 802\] с помощью БПФ в поле вычетов $\mathbb{Z}/p\mathbb{Z}$. Обосновать выбор параметров алгоритма.
		\item Вычислить произведение в поле Ферма: \[152 \cdot 194\ (\textrm{mod}\ 257).\]
		\item Возвести целое число в степень с помощью $4$-арного алгоритма: \[2^{114}.\]
		\item С помощью теста Соловея-Штрассена определить, является ли простым число $N = 439$. Требуется дать ответ с вероятностью ошибки не~более~$\frac14$.
		\item Определить порядок элемента $571$ поля вычетов $\mathbb{Z}/911\mathbb{Z}$ с помощью построения дерева произведений.
		\item Восстановить параметры $a$, $c$ и $m$ Линейного конгруэнтного генератора \[x_i \equiv a \cdot x_{i-1} + c\ (\textrm{mod}\ m)\] по последовательности подряд идущих выходных значений генератора: \[25, 90, 53, 33, 141, 123.\]
	\end{enumerate}
\clearpage	\subsection{Вариант 2}
	\begin{enumerate}
		\item Вычислить произведение целых чисел \[553986280 \cdot 346672728\] с помощью метода Тоома-Кука. Использовать полином 2-й степени и десятичную систему счисления.
		\item Вычислить произведение целых чисел \[316 \cdot 635\] с помощью БПФ в поле вычетов $\mathbb{Z}/p\mathbb{Z}$. Обосновать выбор параметров алгоритма.
		\item Вычислить произведение в поле Ферма: \[221 \cdot 23\ (\textrm{mod}\ 257).\]
		\item Возвести целое число в степень с помощью бинарного алгоритма: \[3^{82}.\]
		\item С помощью теста Миллера-Рабина определить, является ли простым число $N = 401$. Требуется дать ответ с вероятностью ошибки не~более~$\frac{1}{16}$.
		\item Определить порядок элемента $264$ поля вычетов $\mathbb{Z}/547\mathbb{Z}$ с помощью построения дерева произведений.
		\item Восстановить параметры $a$, $c$ и $m$ Линейного конгруэнтного генератора \[x_i \equiv a \cdot x_{i-1} + c\ (\textrm{mod}\ m)\] по последовательности подряд идущих выходных значений генератора: \[100, 35, 14, 81, 106, 47.\]
	\end{enumerate}
\clearpage	\subsection{Вариант 3}
	\begin{enumerate}
		\item Вычислить произведение целых чисел \[342772644 \cdot 142092061\] с помощью метода Тоома-Кука. Использовать полином 2-й степени и десятичную систему счисления.
		\item Вычислить произведение целых чисел \[129 \cdot 735\] с помощью БПФ в поле вычетов $\mathbb{Z}/p\mathbb{Z}$. Обосновать выбор параметров алгоритма.
		\item Вычислить произведение в поле Ферма: \[150 \cdot 154\ (\textrm{mod}\ 257).\]
		\item Возвести целое число в степень с помощью бинарного алгоритма: \[3^{82}.\]
		\item С помощью теста Соловея-Штрассена определить, является ли простым число $N = 541$. Требуется дать ответ с вероятностью ошибки не~более~$\frac14$.
		\item Определить порядок элемента $631$ поля вычетов $\mathbb{Z}/967\mathbb{Z}$ с помощью построения дерева произведений.
		\item Восстановить параметры $a$, $c$ и $m$ Линейного конгруэнтного генератора \[x_i \equiv a \cdot x_{i-1} + c\ (\textrm{mod}\ m)\] по последовательности подряд идущих выходных значений генератора: \[78, 23, 73, 54, 36, 70.\]
	\end{enumerate}
\clearpage	\subsection{Вариант 4}
	\begin{enumerate}
		\item Вычислить произведение целых чисел \[40175657 \cdot 48824181\] с помощью метода Карацубы. Применять метод рекурсивно, база рекурсии — двузначные числа, которые можно умножать любым способом.\\ Использовать десятичную систему счисления.
		\item Вычислить произведение целых чисел \[92 \cdot 786\] с помощью БПФ в поле вычетов $\mathbb{Z}/p\mathbb{Z}$. Обосновать выбор параметров алгоритма.
		\item Вычислить произведение в поле Мерсенна: \[190 \cdot 118\ (\textrm{mod}\ 255).\]
		\item Возвести целое число в степень с помощью бинарного алгоритма: \[2^{78}.\]
		\item С помощью теста Миллера-Рабина определить, является ли простым число $N = 641$. Требуется дать ответ с вероятностью ошибки не~более~$\frac{1}{16}$.
		\item Определить порядок элемента $127$ поля вычетов $\mathbb{Z}/547\mathbb{Z}$ с помощью построения дерева произведений.
		\item Открытый текст $M$ зашифрован криптосистемой RSA с помощью пары открытых ключей $(e_1, N)$ и $(e_2, N)$:  \[\begin{aligned} C_1 \equiv M^{e_1}\ (\textrm{mod}\ N), \\ C_2 \equiv M^{e_2}\ (\textrm{mod}\ N). \end{aligned}\] Зная пару шифротекстов и открытые ключи: \[C_1 = 17058, C_2 = 8040, e_1 = 11, e_2 = 17, N = 23489,\] восстановить открытый текст $M$ без разложения модуля $N$ на множители.
	\end{enumerate}
\clearpage	\subsection{Вариант 5}
	\begin{enumerate}
		\item Вычислить произведение целых чисел \[50070681 \cdot 144949975\] с помощью метода Тоома-Кука. Использовать полином 2-й степени и десятичную систему счисления.
		\item Вычислить произведение целых чисел \[568 \cdot 395\] с помощью БПФ в поле вычетов $\mathbb{Z}/p\mathbb{Z}$. Обосновать выбор параметров алгоритма.
		\item Вычислить произведение в поле Ферма: \[159 \cdot 172\ (\textrm{mod}\ 257).\]
		\item Возвести целое число в степень с помощью бинарного алгоритма: \[3^{45}.\]
		\item С помощью теста Соловея-Штрассена определить, является ли простым число $N = 461$. Требуется дать ответ с вероятностью ошибки не~более~$\frac14$.
		\item Определить порядок элемента $399$ поля вычетов $\mathbb{Z}/463\mathbb{Z}$ с помощью построения дерева произведений.
		\item Восстановить параметры $a$, $c$ и $m$ Линейного конгруэнтного генератора \[x_i \equiv a \cdot x_{i-1} + c\ (\textrm{mod}\ m)\] по последовательности подряд идущих выходных значений генератора: \[65, 206, 29, 275, 363, 70.\]
	\end{enumerate}
\clearpage
\end{document}
